% !TEX root = ./index.tex

\maketitle
\newpage

\section*{Answers}

\begin{enumerate}

        % #TODO 1.
    \item

        \bigskip

        \begin{enumerate}

                % #TODO 1.a.
            \item

                \bigskip

                % #TODO 1.b.
            \item

                \bigskip

                % #TODO 1.c.
            \item

                \bigskip

                \begin{enumerate}

                        % #TODO 1.c.i.
                    \item

                        \bigskip

                        % #TODO 1.c.ii.
                    \item

                \end{enumerate}

                \bigskip

                % #TODO 1.d.
            \item

                \bigskip

                \begin{enumerate}

                        % #TODO 1.d.i.
                    \item

                        % #TODO 1.d.ii.
                    \item

                        % #TODO 1.d.iii.
                    \item

                        % #TODO 1.d.iiii.
                    \item

                \end{enumerate}

        \end{enumerate}

        \bigskip

        % 2.
    \item Let $G(x) = \int_0^{x^2} g(t) \, dt$ for all $x \in \mathbb{R}$, where $g(t) = e^{\sin t} \sin t$ for all $t \in \mathbb{R}$.

        \bigskip

        \begin{enumerate}

                % 2.a.
            \item Consider when $x \in (0, \infty)$.

                \medskip

                \begin{proposition}
                    \setlength{\baselineskip}{20pt}
                    The minima of $G(x)$ are $x = \sqrt{2 \pi}, \sqrt{4 \pi}, \ldots$, and the maxima of $G(x)$ are $x = \sqrt{\pi}, \sqrt{3 \pi}, \ldots$.
                \end{proposition}

                \smallskip

                \begin{proof}
                    The local extrema points of $G(x)$ are characterised by the relation $G'(x) = 0$, and by the Fundamental Theorem of Calculus, this relation is equivalent to
                    \[ 2 x \left(g(x^2) - g(0)\right) = 0, \]
                    or
                    \[ g(x^2) = g(0). \]
                    Since $g(0) = 0$, the local extrema of $G(x)$ occur at all $x$ such that $g(x^2) = 0$.

                    The function $g(t)$ can be divided into the expressions $e^{\sin t}$ and $\sin t$.
                    Note that $g(t) = 0$ if and only if either of these expressions equal zero.
                    The exponential function is always positive and thus can never equal zero, which means that $g(t) = 0$ if and only if $\sin t = 0$.
                    This occurs at the values of $t$ such that $t = k \pi$ for $k \in \mathbb{Z}$.

                    It follows that $x$ is a local extremum of $G(x)$ if and only if $g(x^2) = 0$, which is when
                    \[ x \in \{\sqrt{\pi}, \sqrt{2 \pi}, \ldots, \sqrt{n \pi}\}, \quad n \in \mathbb{N}. \]

                    To classify these extrema, we have
                    \[ G'(x) = 2 x \left(g(x^2) - g(0)\right) = 2 x \, g(x^2), \]
                    so that
                    \begin{spreadlines}{\baselineskip}
                        \begin{align*}
                            G''(x) &= 2 \, g(x^2) + 4 x^2 \, g'(x^2), \\
                            &= 2 \, e^{\sin x^2} \sin x^2 + 4 x^2 \, \frac{d}{d t} \big|_{x^2} \; e^{\sin t} \sin t, \\
                            &= 2 \, e^{\sin x^2} \sin x^2 + 4 x^2 \, \left[e^{\sin t} (\sin t \cos t + \cos t)\right]_{x^2}, \\
                            &= 2 \, e^{\sin x^2} \sin x^2 + 4 x^2 \, e^{\sin x^2} (\sin x^2 \cos x^2 + \cos x^2), \\
                            &= e^{\sin x^2} \left(2 \, \sin x^2 + 4 x^2 \, (\sin x^2 + 1) \cos x^2\right).
                        \end{align*}
                    \end{spreadlines}

                    When $x = \sqrt{n \pi}$ for some $n \in \mathbb{N}$, we have $\sin x^2 = 0$ and $\cos x^2 = (-1)^n$.

                    The value of $G''(x)$ at an extremum is thus
                    \[
                        G''(x) = e^0 \left(2 \cdot 0 + 4 n \pi (0 + 1) \cdot (-1)^n\right) = 4 n \pi (-1)^n.
                    \]
                    This expression is positive when $n$ is even and negative when $n$ is odd, hence the minima of $G(x)$ are
                    \[ x = \sqrt{2 \pi}, \sqrt{4 \pi}, \ldots, \]
                    and the maxima of $G(x)$ are
                    \[ x = \sqrt{\pi}, \sqrt{3 \pi}, \ldots. \]
                \end{proof}

                \bigskip

                % 2.b.
            \item To find the second order Taylor polynomial $T_2(t)$ of $g(t)$ about $t = 0$, first take successive derivatives of $g(t)$:
                \begin{align*}
                    g(t) &= e^{\sin t} \sin t, \\
                    \intertext{so}
                    g'(t) &= e^{\sin t} \sin t \cos t + e^{\sin t} \cos t, \\
                    &= \left(g(t) + e^{\sin t}\right) \cos t, \\
                    \intertext{and finally,}
                    g''(t) &= \left(g'(t) + e^{\sin t} \cos t\right) \cos t - \left(g(t) + e^{\sin t}\right) \sin t.
                \end{align*}

                Now set $t = 0$ to get
                \begin{align*}
                    g(0) &= e^{0} \cdot 0 = 0, \\
                    g'(0) &= \left(g(0) + e^{0}\right) \cdot 1, \\
                    &= 0 + 1, \\
                    &= 1, \\
                    g''(0) &= \left(g'(0) + e^{0} \cdot 1\right) \cdot 1 - \left(g(0) + e^{0}\right) \cdot 0, \\
                    &= (1 + 1) - 0, \\
                    &= 2.
                \end{align*}

                Substituting these values into the expression for $T_2(t)$ gives
                \begin{align*}
                    T_2(t) &= \frac{1}{0!} \, g(0) \, t^0 + \frac{1}{1!} \, g'(0) \, t^1 + \frac{1} {2!} \, g''(0) \, t^2, \\
                    &= g(0) + g'(0) \, t + \frac{1}{2} \, g''(0) \, t^2, \\
                    &= 0 + 1 t + \frac{2}{2} t^2, \\
                    &= t + t^2.
                \end{align*}

                \bigskip

                % 2.c.
            \item Since $T_2(t)$ is the Taylor polynomial of $g(t)$ and $G(x) = \int_0^{x^2} g(t) \, dt$, consider the function $P_6(x) = \int_0^{x^2} T_2(t) \, dt$.
                From the expression of $T_2(t)$ found in part (b), we have
                \begin{spreadlines}{0.5\baselineskip}
                    \begin{align*}
                        P_6(x) &= \int_0^{x^2} t + t^2 \, dt, \\
                        &= \left[\frac{1}{2} t^2 + \frac{1}{3} t^3\right]_{x^2}, \\
                        &= \frac{1}{2} x^4 + \frac{1}{3} x^6.
                    \end{align*}
                \end{spreadlines}

                \medskip

                \begin{proposition} \label{prop:3c}
                    $P_6(x)$ is the sixth order Taylor polynomial of $G(x)$ about $x = 0$.
                \end{proposition}

                \smallskip

                \begin{proof}
                    We will apply the uniqueness property of \emph{Taylor's theorem} \parencite{taylors-theorem}:
                    \begin{theorem} \label{thm:taylor}
                        Let $f: (a, b) \to \mathbb{R}$ be $k$ times differentiable at $x_0 \in (a, b)$ and $P_k(x)$ be the $k$-th order Taylor polynomial of $f(x)$ about $x = x_0$.
                        Then, there exists a real-valued function $R(x)$ such that
                        \[ f(x) = T(x) + R(x) \]
                        and
                        \[ \lim_{x \to x_0} \; \frac{R(x)}{(x - x_0)^k} = 0. \]
                    \end{theorem}

                    \smallskip

                    It can be proven that the remainder term $R(x)$ is unique to $f(x)$.
                    That is, if for some function $f: (a, b) \to \mathbb{R}$ and a polynomial $P(x)$, the function $R(x) := f(x) - P(x)$ satisfies the limit above, then $P(x)$ is the $k$-th order Taylor polynomial of $f(x)$ about $x = x_0$.

                    Recall from part (b) that $T_2(t) = t + t^2$ was shown to be the second order Taylor polynomial of $g(t)$.
                    Applying \cref{thm:taylor} gives
                    \begin{equation} \label{eq:taylor}
                        \lim_{t \to 0} \; \frac{g(t) - (t + t^2)}{t^2} = 0.
                    \end{equation}

                    We will use this result to show that $\lim_{x \to 0} \; \frac{G(x) - P_6(x)}{x^6} = 0$, with the proof immediately following from the uniqueness of the remainder term.

                    Since $g(t)$ is continuous at $t = 0$, the antiderivative is finite, thus $G(0) = \int_0^0 g(t) \, dt = 0$.
                    This means that the quotient $\frac{G(x) - P_6(x)}{x^6}$ is indeterminant as $x \to 0$.
                    Also, $G(x)$ is differentiable at $x = 0$, hence we can use L'Hospital's rule on the limit:
                    \begin{spreadlines}{0.5\baselineskip}
                        \begin{align*}
                            \lim_{x \to 0} \; \frac{G(x) - P_6(x)}{x^6} &= \lim_{x \to 0} \; \frac{G'(x) - P_6'(x)}{6 x^5}, \\
                            &= \lim_{x \to 0} \; \frac{2 x \, g(x^2) - (2 x^3 + 2 x^5)}{6 x^5}, \\
                            &= \lim_{x \to 0} \; \frac{g(x^2) - (x^2 + x^4)}{3 x^4}.
                        \end{align*}
                    \end{spreadlines}

                    Now, $\lim_{t \to 0} \frac{g(t) - (t + t^2)}{t^2} = 0$ from \cref{eq:taylor}, and $\lim_{x \to 0} \; x^2 = 0$.
                    Hence by the substitution law of limits,
                    \[ \lim_{x \to 0} \; \frac{g(x^2) - (x^2 + x^4)}{3 x^4} = \frac{0}{3} = 0. \]

                    We have thus proven that $\lim_{x \to 0} \; \frac{G(x) - P_6(x)}{x^6} = 0$.
                    Because the remainder term $R(x) = G(x) - P_6(x)$ is the unique function which statisfies the limit $\lim_{x \to 0} \; \frac{R(x)}{x^6} = 0$, the polynomial $P_6(x)$ is the sixth order Taylor polynomial of $G(x)$ by \cref{thm:taylor}.
                \end{proof}

                \bigskip

                % 2.d.
            \item The result of \cref{prop:3c} can be written as
                \[ \lim_{x \to 0} \; \frac{G(x) - (\frac{1}{2} x^4 + \frac{1}{3} x^6)}{x^6} = 0, \]
                or
                \begin{equation} \label{eq:taylor-1}
                    \lim_{x \to 0} \;\frac{G(x) - \frac{1}{2} x^4}{x^6} = \frac{1}{3}.
                \end{equation}
                We will use this equation to prove the following proposition.

                \medskip

                \begin{proposition}
                    $\lim_{x \to 0} \; \frac{G(x) - 2 \int_0^x \sin(t^3) \, dt}{x^6} = \frac{1}{3}$.
                \end{proposition}

                \smallskip

                \begin{proof}
                    Consider the fifth order Taylor polynomial of $h(t) := \sin(t^3)$ about $t = 0$. Using the substitution method on the Taylor polynomial of $t \to \sin t$, this is given by
                    \[ h(t) \approx t^3. \]
                    Taking the primitive yields the sixth order Taylor polynomial of $H(x) := \int_0^x h(t) \, dt$ about $x = 0$:
                    \[ H(x) \approx \frac{1}{4} x^4. \]
                    Thus, by \cref{thm:taylor},
                    \begin{equation} \label{eq:taylor-2}
                        \lim_{x \to 0} \; \frac{H(x) - \frac{1}{4}{x^4}}{x^6} = 0.
                    \end{equation}

                    To complete the proof, simply evaluate $(\ref{eq:taylor-1}) - 2 \, (\ref{eq:taylor-2})$:
                    \begin{spreadlines}{\baselineskip}
                        \begin{align*}
                            \text{LHS} &= \lim_{x \to 0} \; \frac{G(x) - \frac{1}{2} x^4}{x^6} - \lim_{x \to 0} \; \frac{2 \, H(x) - \frac{1}{2}{x^4}}{x^6}, \\
                            &= \lim_{x \to 0} \; \frac{\left(G(x) - \frac{1}{2} x^4\right) - \left(2 \, H(x) - \frac{1}{2} x^4\right)}{x^6}, \\
                            &= \lim_{x \to 0} \; \frac{G(x) - 2 H(x)}{x^6}, \\
                            \text{RHS} &= \frac{1}{3} - 2 \cdot 0 = \frac{1}{3}.
                        \end{align*}
                    \end{spreadlines}

                    Ergo, $\lim_{x \to 0} \; \frac{G(x) - 2 \int_0^x \sin(t^3) \, dt}{x^6} = \frac{1}{3}$.
                \end{proof}

        \end{enumerate}

        \bigskip

        % 3.
    \item Let $M_n(\mathbb{R})$ denote the set of $n \times n$ matrices with entries in $\mathbb{R}$.
        Define the \emph{trace} function $\text{tr}: M_n(\mathbb{R}) \to \mathbb{R}$ as
        \[ \text{tr}(A) := \sum_{i = 1}^n a_{ii}, \]
        where $a_{ii}$ is taken to be the term in the $i$-th row and $i$-th column of the matrix $A$.

        The dot product on $M_n(\mathbb{R})$ is defined as
        \[ A \cdot B := \text{tr}(A^T B). \]

        \bigskip

        \begin{enumerate}

                % 3.a
            \item

                % 3.b
            \item

                % 3.c
            \item

        \end{enumerate}

        \bigskip

        % #TODO 4.
    \item
        \begin{definition} \label{def:invariant-subspace}
            Suppose that $T: \mathbb{R}^n \to \mathbb{R}^n$ is a linear transformation.
            A subspace $S$ of $\mathbb{R}^n$ is an \emph{invariant subspace} of $T$ if $T(\vb{x}) \in S$ for all $\vb{x} \in S$.
        \end{definition}

        \bigskip

        \begin{enumerate}

                % #TODO 4.a.
            \item
                \begin{proposition}
                    Every eigenspace $E_{\lambda}(A_T)$ is an invariant subspace of $T$.
                \end{proposition}

                \smallskip

                \begin{proof}

                \end{proof}

                \bigskip

                % #TODO 4.b.
            \item

        \end{enumerate}

        \bigskip

        % #TODO 5.
    \item

\end{enumerate}
